\documentclass[10pt,a4paper]{amsart}
\usepackage[margin=19mm]{geometry}
\usepackage{amsmath,amssymb,mathtools}
\usepackage[scr=rsfs,cal=euler]{mathalfa}
\usepackage{xfrac}
\usepackage[unicode,hidelinks,bookmarks=false]{hyperref}
\newcommand{\ie}{i.e.\ }
\newcommand{\cf}{cf.\ }
\newcommand{\resp}{respectively\ }
\newcommand{\Subsection}{\subsection}
\providecommand{\nobreakdash}{\nobreak\hbox{-}}
\setlength{\emergencystretch}{2em}
\multlinegap0pt
\usepackage{bm}
\usepackage{bbm}
\usepackage{dsfont}
\usepackage{xspace}
\usepackage{cancel}
\usepackage{mathtools}
\usepackage{amsthm}
\makeatletter
\@ifundefined{lemma}{\newtheorem{lemma}{Lemma}}{}
\@ifundefined{proposition}{\newtheorem{proposition}{Proposition}}{}
\@ifundefined{theorem}{\newtheorem{theorem}{Theorem}}{}
\makeatother
\title[Taylor polynomials on left-quotients of Carnot groups]{Taylor polynomials on left-quotients of Carnot groups}
\author{Alessandro Ottazzi}
\date{Source: arXiv:2512.12239v3 (CC BY 4.0)}
\begin{document}
\maketitle

\noindent\emph{Source and licence.} Taylor polynomials on left-quotients of Carnot groups. Version arXiv:2512.12239v3 (CC BY 4.0), distributed under the Creative Commons Attribution 4.0 International licence, as recorded in the arXiv licence metadata for this exact version and shown on the abstract page https://arxiv.org/abs/2512.12239v3. Full rights notice: licenses/arxiv-2512.12239-CC-BY-4.0.txt.\par\smallskip

\noindent\textit{Editorial scope.} Counted unit: the Theorem whose source label is \texttt{TaylorQuotients} in the article (source0.tex lines 830--850, source label \texttt{TaylorQuotients}), delivered together with the article's own complete original proof of it (source0.tex lines 852--928, one \texttt{proof} environment of 77 lines). No mathematics is written, repaired or re-derived: statement and proof are the article's own text line for line. Required context retained uncounted: submetrylift (lines 308--313); Taylor (lines 477--487); prop:compatibility (lines 563--577); lem:sup-transfer (lines 645--669). Every definition, display and label of those lines is retained unchanged. Omitted from this package and not redistributed: 1--307, 314--476, 488--562, 578--644, 670--829, 851--851, 929--1566. Nothing inside a retained excerpt is omitted or altered: every retained body is delivered exactly as the article's lines are written, so no omission receipt of any kind accompanies this package. Citations: the retained excerpts carry no citation command at all, so no bibliography entry of the article is transcribed and no reference is redistributed. Reference adaptation: none was needed and none was made; every reference inside the delivered unit resolves inside it, so no unresolved reference or citation is produced. Printed slips kept literal: no printed word, symbol, hypothesis or number of the counted statement or of its proof was altered. Layout and class: the article's own class is \texttt{amsart} with packages including mathrsfs, hyperref, amsmath, color, enumerate, amssymb, amsfonts, graphicx; the wrapper is the lane's pinned \texttt{amsart} editorial wrapper at 10pt on A4, which restates the theorem-like environments the retained text needs and adds the article's own macro definitions that the retained text needs (none), with extra wrapper packages (bm, bbm, dsfont, xspace, cancel, mathtools). The delivered numbering is the unit's own. Assets: no retained span contains a figure environment, a graphics-inclusion command, a PDF-page inclusion, a table or a TikZ picture; no image file is copied into this package, the package declares no asset dependency and no image file is redistributed with it. Licence: version arXiv:2512.12239v3 of this article is distributed under the Creative Commons Attribution 4.0 International licence, as recorded in the arXiv licence metadata for this exact version and shown on the abstract page https://arxiv.org/abs/2512.12239v3. A full rights notice is delivered at licenses/arxiv-2512.12239-CC-BY-4.0.txt. Characters: every retained excerpt is pure ASCII, so the delivered engine, XeLaTeX with its default Latin Modern text font, reports no missing character.
\begin{proposition}\label{submetrylift}
    For every $p,q\in M$ and $g_0\in G$ so that $\Pi(g_0)=q$, there is $g\in G$ with $\Pi(g)=p$ so that
  $$
  d_M(p,q) = d_G(g,g_0).
  $$
\end{proposition}

\begin{theorem}[Taylor theorem]\label{Taylor}
    For every positive integer $k$ there is a constant $C_k$ such that for all $F\in C^k_{\mathrm{hor}}(G)$ and all $g,g_0\in G$,
    $$
    |F(g g_0)-\tilde P_k(F,g)(g g_0)| \leq C_k d_G(g_0,0)^k
    \tilde \eta_G(g,b^k d_G(g_0,0)),
    $$
where $b$ is as in the previous theorem and 
$$
\tilde \eta_G(g,b^k d_G(g_0,0)) = \sup\left\{|\tilde X_{\mathrm{hor}}^IF(gu)-\tilde X_{\mathrm{hor}}^IF(g)|\,:\, d(I)=k, \, d_G(u,0)\leq b^kd_G(g_0,0)\right\}.
$$   
\end{theorem}

\begin{proposition}
\label{prop:compatibility}
Let \(f\in C_{\mathrm{hor}}^{k}(M)\) and \(\tilde f := f\circ\Pi\).
Fix \(q\in M\) and choose any \(g_0\in G\) with \(\Pi(g_0)=q\).
Let \(\tilde P_k(\tilde f,g_0)\) denote the  Taylor polynomial of homogeneous degree \(k\) for \(\tilde f\) at \(g_0\).
Then:
\begin{enumerate}
\item \(\tilde P_k(\tilde f,g_0)\) is left \(H\)-invariant and it defines a unique polynomial \(\bar P_k\) on \(H\backslash G\).
\item The restriction of \(\bar P_k\) to \(M\) (via the identification \(\psi\)) coincides with the Taylor polynomial on \(M\), namely
\[
P_{k}(f,q) \;=\; \bar P_k\circ\psi \;=\; \tilde P_k(\tilde f,g_0)\big|_{M}.
\]
\item The polynomial \(\tilde P_k(\tilde f,g_0)\) (and hence \(\bar P_k\) and \(P_k(f,q)\)) is independent of the choice of representative \(g_0\) with \(\Pi(g_0)=q\).
\end{enumerate}
\end{proposition}

\begin{lemma}
\label{lem:sup-transfer}
% Let $G$ be a Carnot group, $H\subset G$ a homogeneous subgroup, and
% $\pi:G\to H\backslash G$ the left quotient map.
% Let $M\subset G$ be a global cross--section, and set
% \[
% \Pi := (\pi|_M)^{-1}\circ \pi : G\to M.
% \]

Let $\Phi:G\to\mathbb R$ be a left $H$-invariant function, and define
$\phi:M\to\mathbb R$ by
\[
\phi := \Phi\circ \iota,
\]
where $\iota:M\hookrightarrow G$ is the inclusion.
% (equivalently, $\phi(m)=\Phi(g)$ for any
% $g\in G$ with $\Pi(g)=m$).
Then for every $q\in M$ and every $r>0$,
\begin{equation}\label{eq:sup-transfer}
\sup \left\{|\phi(p)|\,:\, p\in M, d_M(p,q)\le r\right\}
=
\sup \left\{|\Phi(g)|\,:\,g\in G, d_G(g,g_0)\le r\right\},
\end{equation}
where $g_0\in G$ is any point such that $\Pi(g_0)=q$.
\end{lemma}

\begin{theorem}[Taylor Theorem for $M$]\label{TaylorQuotients}
For every positive integer $k$ there exist constants $C_k>0$ and $b>0$ such that
for all $f\in C^k_{\mathrm{hor}}(M)$ and all $p,q\in M$,
\begin{equation}\label{eq:Taylor-M}
\big| f(p)- P_k(f,q)(p)\big|
\leq
C_k\, d_M(p,q)^k\,
\eta_M\big(q,b^k d_M(p,q)\big),
\end{equation}
where $P_k(f,q)$ is the  Taylor polynomial of degree $k$ of $f$ at $q$,
and
\[
\eta_M(q,r)
:=
\sup\big\{
\big| X_{\mathrm{hor}}^I f(n)-X_{\mathrm{hor}}^I f(q)\big|
\;\big|\;
d(I)=k,\ d_M(n,q)\le r
\big\}.
\]
\end{theorem}

\begin{proof}
Fix $f\in C^k_{\mathrm{hor}}(M)$ and $q\in M$.  Define the left $H$-invariant lift
\(
\tilde f = f\circ \Pi \in C^k_{\mathrm{hor}}(G).
\)
Choose any $g_0\in G$ with $\Pi(g_0)=q$.
Let $\tilde P_k(\tilde f,g_0)$ be the  Taylor polynomial of degree $k$ for $\tilde f$
centered at $g_0$, and let $P_k(f,q)$ be the Taylor polynomial on $M$.
By Proposition~\ref{prop:compatibility}, \(
P_k(f,q) = \tilde P_k(\tilde f,g_0)\big|_{M}.
\)
Now fix $p\in M$. 
By Proposition~\ref{submetrylift}, there exists
$g\in G$ such that
\(
\Pi(g)=p\) and \( d_G(g,g_0)=d_M(p,q).
\)
Then
\[
f(p)-P_k(f,q)(p)
=
\tilde f(g) - \tilde P_k(\tilde f,g_0)(g).
\]
Applying Theorem~\ref{Taylor} to $\tilde f$ and the pair $(g,g_0)$ yields
\begin{equation}\label{eq:TQ-step1}
\big| f(p)-P_k(f,q)(p)\big|
\leq
C_k\, d_G(g,g_0)^k\, \tilde \eta_G\big(g_0,b^k d_G(g,g_0)\big).
\end{equation}
Since $d_G(g,g_0)=d_M(p,q)$ by construction, the factor in front is exactly
$C_k d_M(p,q)^k$.  It remains to rewrite the quantity $\tilde \eta_G$ in terms of $\eta_M$. This is done by applying  Lemma~\ref{lem:sup-transfer} to $\phi(n) = X_{\mathrm{hor}}^I f(n)-X_{\mathrm{hor}}^I f(q)$ and $\Phi(\cdot) = \tilde X_{\mathrm{hor}}^I \tilde f(\cdot)-\tilde X_{\mathrm{hor}}^I \tilde f(g_0)$.

% Let $I$ be any multi-index with $d(I)=k$, and let $n\in G$.
% Set $m := \Pi(n)\in M$.  By definition of $\Pi$ we have $\pi(n)=\pi(m)$, and by construction
% of the vector fields $X_j$ from $\tilde X_j$ via the quotient, we have
% \[
% \tilde X_{\mathrm{hor}}^I \tilde f(n)
% = \tilde X_{\mathrm{hor}}^I (f\circ \Pi)(n)
% = X_{\mathrm{hor}}^I f(m),
% \]
% and similarly
% \[
% \tilde X_{\mathrm{hor}}^I \tilde f(g_0) = X_{\mathrm{hor}}^I f(q).
% \]
% Moreover, since $\Pi$ is a submetry, the quotient distance is
% $1$--Lipschitz with respect to $d_G$:
% \[
% d_M(m,q) \le d_G(n,g_0) \quad \text{for all } n\in G.
% \]
% In particular, if $d_G(n,g_0)\le r$ then $d_M(m,q)\le r$.
% Therefore, for any $r>0$,
% \begin{align*}
% \sup 
% \left\{\big|\tilde X_{\mathrm{hor}}^I \tilde f(n)-\tilde X_{\mathrm{hor}}^I \tilde f(g_0)\big|\,:\,\right. &\left.
% d(I)=k, d_G(n,g_0)\le r
% \right\}
% \\
% &=
% \sup
% \left\{\big|X_{\mathrm{hor}}^I f(m)-X_{\mathrm{hor}}^I f(q)\big|\,:\, 
% d(I)=k, d_M(m,q)\le r
% \right\}.
% \end{align*}
% That is,
% \[
% \eta_G(g_0,r) = \eta_M(q,r) \quad \text{for all } r>0.
% \]

% Substituting $r=b^k d_G(g,g_0)=b^k d_M(p,q)$ into~\eqref{eq:TQ-step1} and using
% $\eta_G(g_0,r)=\eta_M(q,r)$ gives
% \[
% \big| f(p)-P_k(f,q)(p)\big|
% \leq
% C_k\, d_M(p,q)^k\, \eta_M\big(q,b^k d_M(p,q)\big),
% \]
% which is exactly \eqref{eq:Taylor-M}.  
\end{proof}

\end{document}
